\documentclass[../../main.tex]{subfiles}

\begin{document}

\chapter{Landasan Matematika untuk Kriptografi}
\label{chap:landasan-matematika}

\begin{learningoutcome}
Setelah mempelajari bab ini, mahasiswa diharapkan mampu:
\begin{itemize}
    \item Mengaplikasikan konsep aritmetika modulo, kekongruenan, dan relatif prima dalam operasi kriptografi (Sub-CPMK 2.3).
    \item Menentukan balikan (invers) modulo menggunakan algoritma Euclidean (Sub-CPMK 2.3).
    \item Menjelaskan fungsi totient Euler, Teorema Euler, akar primitif, dan persoalan logaritma diskrit serta kaitannya dengan algoritma kunci publik (Sub-CPMK 2.3).
    \item Menjelaskan konsep dasar teori informasi, khususnya mengenai entropi Shannon, dan kaitannya dengan keamanan cipherteks (Sub-CPMK 2.3).
\end{itemize}
\end{learningoutcome}

\subfile{section-01}
\subfile{section-02}
\section{Contoh Terhitung}

\begin{example}[Algoritma Euclidean dan invers modulo]
\begin{enumerate}
    \item \textbf{PBB dan koefisien B\'ezout dari pasangan $(1970,1066)$.} Terapkan algoritma pembagian berulang berikut untuk memperoleh sisa terakhirnya:
    \[
    \begin{aligned}
    1970 &= 1 \cdot 1066 + 904\\
    1066 &= 1 \cdot 904 + 162\\
    904  &= 5 \cdot 162 + 94\\
    162  &= 1 \cdot 94 + 68\\
    94   &= 1 \cdot 68 + 26\\
    68   &= 2 \cdot 26 + 16\\
    26   &= 1 \cdot 16 + 10\\
    16   &= 1 \cdot 10 + 6\\
    10   &= 1 \cdot 6 + 4\\
    6    &= 1 \cdot 4 + 2\\
    4    &= 2 \cdot 2 + 0
    \end{aligned}
    \]
    Sisa bukan nol terakhir adalah $2$, jadi $\text{PBB}(1970,1066) = 2$. Substitusi balik memberi koefisien B\'ezout
    \[ 1970 \cdot (-204) + 1066 \cdot 377 = 2. \]
    \item \textbf{Invers modulo $550 \bmod 1769$.} Karena $\text{PBB}(550,1769) = 1$, inversnya ada. Uji nilai $x = 550$:
    \[ 550 \cdot 550 = 302\,500 = 171 \cdot 1769 + 1, \]
    sehingga $550 \cdot 550 \equiv 1 \pmod{1769}$, yaitu $550^{-1} \equiv 550 \pmod{1769}$.
\end{enumerate}
\end{example}

\begin{example}[Entropi Shannon sebuah distribusi]
Pesan $X = \code{AABBCBDB}$ terdiri atas 8 simbol dengan peluang $p(A)=2/8$, $p(B)=4/8$, $p(C)=1/8$, dan $p(D)=1/8$. Substitusikan ke rumus entropi:
\[
\begin{aligned}
H(X) &= -\sum_{i} p(x_i) \log_2 p(x_i)\\
     &= -\left(\tfrac14 \log_2 \tfrac14 + \tfrac12 \log_2 \tfrac12
        + \tfrac18 \log_2 \tfrac18 + \tfrac18 \log_2 \tfrac18\right)\\
     &= -\left(\tfrac14(-2) + \tfrac12(-1) + \tfrac18(-3) + \tfrac18(-3)\right)\\
     &= 1{,}75\ \text{bit/simbol}.
\end{aligned}
\]
Karena entropi maksimum untuk 4 simbol adalah $\log_2 4 = 2$ bit/simbol, distribusi yang tidak seragam ini menghasilkan entropi yang lebih rendah dan karena itu lebih mudah dianalisis secara statistik.
\end{example}

\begin{example}[Fungsi totient Euler dan Teorema Euler]
Hitung $\varphi(20)$, yaitu banyaknya bilangan pada rentang $1 \le k \le 20$ yang
relatif prima terhadap 20. Bilangan-bilangan itu adalah
\[ 1,\ 3,\ 7,\ 9,\ 11,\ 13,\ 17,\ 19, \]
seluruhnya delapan buah, jadi $\varphi(20) = 8$. Pemeriksaan melalui sifat
multiplikatif memberi hasil yang sama: karena $20 = 4 \cdot 5$ dan
$\text{PBB}(4,5)=1$,
\[ \varphi(20) = \varphi(4) \cdot \varphi(5) = 2 \cdot 4 = 8. \]

Selanjutnya periksa Teorema Euler untuk $a = 3$ dan $n = 10$. Karena
$\text{PBB}(3,10) = 1$ dan $\varphi(10) = \varphi(2)\cdot\varphi(5) = 1 \cdot 4 = 4$,
teorema menyatakan $3^4 \equiv 1 \pmod{10}$. Pemeriksaannya:
\[ 3^4 = 81 = 8 \cdot 10 + 1, \]
sehingga benar bahwa $3^4 \equiv 1 \pmod{10}$. Perhatikan bahwa yang diperiksa bukan
nilai $3^4$ itu sendiri, melainkan sisanya ketika dibagi 10.
\end{example}

\begin{example}[Akar primitif dan logaritma diskrit]
Tinjau modulo $p = 7$. Untuk $g = 3$, hitung pangkat-pangkatnya secara berurutan:
\[
\begin{aligned}
3^1 &= 3 \equiv 3, &\qquad 3^2 &= 9 \equiv 2,\\
3^3 &= 27 \equiv 6, &\qquad 3^4 &= 81 \equiv 4,\\
3^5 &= 243 \equiv 5, &\qquad 3^6 &= 729 \equiv 1.
\end{aligned}
\]
Siklusnya menghasilkan $3, 2, 6, 4, 5, 1$ --- keenam anggota $\{1,2,3,4,5,6\}$ muncul
tepat satu kali, sehingga $g = 3$ adalah akar primitif modulo 7.

Sebagai pembanding, untuk $g = 2$:
\[ 2^1 \equiv 2, \qquad 2^2 \equiv 4, \qquad 2^3 = 8 \equiv 1, \]
sehingga siklusnya hanya berpanjang tiga dan hanya menghasilkan $\{1,2,4\}$. Karena
tidak seluruh anggota $\{1,\dots,6\}$ tercakup, $g = 2$ bukan akar primitif modulo 7.

Selanjutnya tinjau logaritma diskrit $7^x \equiv 15 \pmod{41}$. Hitung pangkat 7
secara berurutan sambil mengurangi modulo 41:
\[
\begin{aligned}
7^1 &= 7 \equiv 7, &\qquad 7^2 &= 49 = 41 + 8 \equiv 8,\\
7^3 &= 7 \cdot 8 = 56 = 41 + 15 \equiv 15.
\end{aligned}
\]
Jadi $x = 3$. Perhatikan bahwa menghitung $7^3 \bmod 41$ hanya menuntut dua perkalian,
sedangkan menemukan $x = 3$ dari $7^x \equiv 15$ pada modulus yang berukuran realistis
--- misalnya 2048 bit --- tidak memiliki cara penyelesaian yang praktis.
\end{example}

\begin{example}[Entropi plainteks berita Kota Ternate]
Plainteks berita Kota Ternate pada Bagian~\ref{sec:teori-informasi} memuat 296 huruf
setelah spasi dan tanda baca dibuang, dengan 21 huruf berbeda. Enam huruf yang paling
sering muncul adalah:

\begin{center}
\begin{tabular}{@{}lcccccc@{}}
\hline
Huruf & A & N & E & I & K & T \\
Cacah & 59 & 34 & 29 & 22 & 17 & 17 \\
Peluang & $0{,}19932$ & $0{,}11486$ & $0{,}09797$ & $0{,}07432$ & $0{,}05743$ & $0{,}05743$ \\
\hline
\end{tabular}
\end{center}

Kontribusi keenam huruf ini terhadap entropi dihitung dengan suku $-p \log_2 p$:
\[
\begin{aligned}
-p(A)\log_2 p(A) &= 0{,}19932 \times 2{,}327 &= 0{,}4638\\
-p(N)\log_2 p(N) &= 0{,}11486 \times 3{,}122 &= 0{,}3586\\
-p(E)\log_2 p(E) &= 0{,}09797 \times 3{,}352 &= 0{,}3284\\
-p(I)\log_2 p(I) &= 0{,}07432 \times 3{,}750 &= 0{,}2787\\
-p(K)\log_2 p(K) &= 0{,}05743 \times 4{,}122 &= 0{,}2367\\
-p(T)\log_2 p(T) &= 0{,}05743 \times 4{,}122 &= 0{,}2367
\end{aligned}
\]
Jumlah keenamnya adalah $1{,}9029$ bit. Lima belas huruf sisanya --- yang masing-masing
muncul jauh lebih jarang --- menyumbang $1{,}9959$ bit lagi, sehingga totalnya
\[ H(P) = 1{,}9029 + 1{,}9959 = 3{,}8988\ \text{bit/karakter}. \]
Dibandingkan dengan entropi maksimum alfabet 26 huruf, pencapaiannya
\[ \frac{3{,}8988}{4{,}7004} \times 100\% \approx 82{,}95\%. \]
Perhatikan bahwa lima belas huruf yang jarang muncul justru menyumbang bagian entropi
yang lebih besar daripada enam huruf tersering: inilah wujud nyata dari tafsir entropi
sebagai ``rata-rata kejutan''. Huruf yang jarang mengagetkan, dan kejutan itulah yang
mengandung informasi.
\end{example}


\begin{obeactivity}
\textbf{Aktivitas 3.1: Eksperimen Modulo dan Invers}
1. Pilihlah dua bilangan acak $a$ dan $n$.
2. Periksa apakah $\text{PBB}(a, n) = 1$.
3. Jika ya, carilah invers modulo $a \pmod{n}$ menggunakan metode coba-coba atau Algoritma Euclidean.
4. Diskusikan apa yang terjadi jika $\text{PBB}(a, n) \neq 1$. Apakah invers modulo tetap ada?
\end{obeactivity}


\begin{obereflection}
\textbf{Latihan 3.1}
\begin{enumerate}
    \item Hitunglah nilai dari $-24 \pmod{11}$.
    \item Tentukan invers modulo dari $12 \pmod{5}$.
    \item Diketahui sebuah pesan memiliki huruf A, B, C, D dengan probabilitas
          masing-masing $p(A)=0{,}25$, $p(B)=0{,}25$, $p(C)=0{,}25$, $p(D)=0{,}25$.
          Hitunglah entropi pesan tersebut.
    \item Tentukan $\text{PBB}(48, 18)$ dengan Algoritma Euclidean, lalu nyatakan
          hasilnya sebagai kombinasi lanjar $m \cdot 48 + n \cdot 18$.
    \item Hitunglah $\varphi(20)$, $\varphi(7)$, dan $\varphi(8)$. Jelaskan mengapa
          hasil $\varphi(20)$ sama dengan $\varphi(4)\cdot\varphi(5)$.
    \item Periksa Teorema Euler untuk $a = 3$ dan $n = 10$: hitung $\varphi(10)$ dan
          tunjukkan bahwa $3^{\varphi(10)} \equiv 1 \pmod{10}$.
    \item Tunjukkan bahwa $g = 3$ adalah akar primitif modulo 7, dan bahwa $g = 2$
          bukan. Tuliskan seluruh pangkatnya sampai kembali ke 1.
    \item Diketahui $7^x \equiv 15 \pmod{41}$. Carilah $x$. Jelaskan mengapa
          menghitung $7^x$ dari $x$ yang diketahui jauh lebih mudah daripada
          menentukan $x$ dari hasilnya.
    \item Sebuah pesan 200 huruf memiliki entropi 4,2 bit/karakter. Berapa persen
          entropi maksimum alfabet 26 huruf yang dicapainya, dan apa artinya bagi
          kekuatan cipherteks tersebut terhadap analisis frekuensi?
\end{enumerate}
\end{obereflection}

\section{Rangkuman}
\begin{itemize}
    \item Aritmetika modulo menghasilkan sisa bagi $a \bmod n$ dengan $a = qn + r$ dan $0 \le r < n$; representasi ini menjadi dasar seluruh operasi kriptografi modern.
    \item Untuk bilangan negatif, hasil $a \bmod n$ tetap harus berada pada rentang $0 \le r < n$, sehingga $-41 \bmod 9 = 4$ dan $-24 \bmod 11 = 9$, bukan bilangan negatif.
    \item Dua bilangan kongruen modulo $n$ jika $n$ membagi habis selisihnya, ditulis $a \equiv b \pmod{n}$. Penjumlahan, pengurangan, dan perkalian dapat dilakukan pada keduanya tanpa mengubah kebenarannya.
    \item Pencoretan pada kekongruenan tidak selalu sah: dari $ac \equiv bc \pmod n$ hanya boleh disimpulkan $a \equiv b \pmod n$ bila $\text{PBB}(c,n) = 1$.
    \item $a$ dan $b$ relatif prima bila $\text{PBB}(a,b) = 1$; syarat ini menentukan keberadaan invers modulo.
    \item Algoritma Euclidean mencari PBB melalui pembagian berulang, dan versi diperluasnya (substitusi balik) menghasilkan koefisien B\'ezout $ma + nb = \text{PBB}(a,b)$ untuk menentukan invers modulo.
    \item Invers modulo $x = a^{-1} \bmod n$ memenuhi $a \cdot x \equiv 1 \pmod{n}$ dan hanya ada jika $\text{PBB}(a,n)=1$.
    \item Bilangan prima adalah bilangan lebih besar dari 1 yang hanya habis dibagi 1 dan dirinya sendiri; Teorema Dasar Aritmetika menjamin pemfaktoran tunggal, sedangkan kesulitan memfaktorkan hasil kali dua prima besar menjadi dasar keamanan RSA.
    \item Fungsi totient Euler $\varphi(n)$ menghitung banyaknya bilangan pada $1 \le k \le n$ yang relatif prima terhadap $n$; berlaku $\varphi(p) = p-1$ dan, untuk $n = pq$ hasil kali dua prima berbeda, $\varphi(n) = (p-1)(q-1)$.
    \item Teorema Euler menyatakan $a^{\varphi(n)} \equiv 1 \pmod n$ bila $\text{PBB}(a,n)=1$, dan Teorema Fermat Kecil $a^{p-1} \equiv 1 \pmod p$ adalah kasus khususnya untuk $p$ prima.
    \item Akar primitif modulo $p$ adalah $g$ yang pangkat-pangkatnya membangkitkan seluruh $\{1,\dots,p-1\}$; kebalikannya, menentukan $x$ dari $g^x \equiv h \pmod p$ adalah persoalan logaritma diskrit yang dianggap sulit dan menjadi dasar Diffie--Hellman serta ElGamal.
    \item Entropi Shannon $H(X) = -\sum_i p(x_i)\log_2 p(x_i)$ mengukur ketidakpastian pesan dalam bit, dengan rentang $0 \le H(X) \le \log_2 n$.
    \item Entropi maksimum alfabet 26 huruf adalah $\log_2 26 = 4{,}7004$ bit/karakter, sedangkan 256 karakter ASCII memberi $\log_2 256 = 8$ bit/karakter.
    \item Teks berbahasa alami memiliki entropi jauh di bawah maksimumnya karena redundansi; selisih inilah yang membuat analisis frekuensi mungkin dilakukan.
    \item Pada contoh berita Kota Ternate, entropi plainteks $H(P) = 3{,}8988$ bit/karakter (82,95\% dari maksimum) naik menjadi $H(C) = 4{,}3035$ bit/karakter (91,55\% dari maksimum) setelah dienkripsi, tampak dari sebaran huruf cipherteks yang jauh lebih rata.
    \item Cipherteks yang baik berentropi tinggi mendekati $\log_2 26 \approx 4{,}70$ bit/karakter sehingga sulit dipecahkan dengan analisis frekuensi; namun entropi tinggi bukan bukti keamanan, melainkan hanya alat diagnosis.
\end{itemize}


\begin{obeassessment}
\textbf{Kuis Bab 3}
\begin{enumerate}
    \item Mengapa pemahaman tentang bilangan prima dan relatif prima sangat penting
          dalam algoritma kunci publik seperti RSA?
    \item Jika sebuah cipherteks memiliki nilai entropi yang sangat rendah, apa
          implikasinya terhadap keamanan pesan tersebut?
    \item Hitung $-37 \bmod 12$ dan $5^{-1} \bmod 12$. Tunjukkan langkah pemeriksaannya.
    \item Hitung $\varphi(21)$ dengan dua cara: mendaftar bilangan yang relatif prima
          terhadap 21, dan memakai sifat multiplikatif karena $21 = 3 \cdot 7$.
    \item Jelaskan mengapa keberadaan invers modulo $a^{-1} \bmod n$ bergantung pada
          nilai $\text{PBB}(a,n)$, dan buktikan dengan Identitas B\'ezout.
    \item Sebuah cipherteks 400 karakter memiliki entropi 4,55 bit/karakter. Berapa
          persen entropi maksimum alfabet 26 huruf yang dicapainya, dan mengapa angka
          itu belum cukup untuk menyimpulkan bahwa ciphernya aman?
    \item Mengapa menentukan $x$ pada $g^x \equiv h \pmod p$ dianggap sulit, padahal
          menghitung $g^x \bmod p$ dari $x$ yang diketahui sangat mudah?
\end{enumerate}
\end{obeassessment}
\begin{competencychecklist}
\begin{tabularx}{\linewidth}{|X|c|}
\hline
Indikator Kompetensi & Tercapai \\
\hline
Saya dapat melakukan operasi aritmetika modulo dan kekongruenan & [ ] \\
\hline
Saya dapat menentukan apakah dua bilangan relatif prima & [ ] \\
\hline
Saya dapat menghitung PBB dengan Algoritma Euclidean dan koefisien B\'ezout-nya & [ ] \\
\hline
Saya dapat menghitung invers modulo & [ ] \\
\hline
Saya dapat menghitung fungsi totient Euler dan menerapkan Teorema Euler & [ ] \\
\hline
Saya dapat menjelaskan akar primitif dan persoalan logaritma diskrit & [ ] \\
\hline
Saya dapat menjelaskan konsep entropi dan menghitung nilainya & [ ] \\
\hline
Saya dapat menafsirkan entropi cipherteks dan kaitannya dengan analisis frekuensi & [ ] \\
\hline
\end{tabularx}
\end{competencychecklist}


\end{document}
